% Chapter 3 — includable by main.tex
\chapter{Konačni automati}
\label{chap:3}

\begin{quote}
\itshape
Regularni izrazi opisuju što jezik sadrži --- ali kako provjeriti
pripada li konkretna riječ nekom jeziku? Za to trebamo stroj.
Konačni automat je najjednostavniji takav stroj: čita riječ znak po
znak i na kraju odlučuje da ili ne, bez ikakve memorije osim
trenutnog stanja. Ovaj jednostavni model ispostavlja se jednako
moćnim kao i regularni izrazi.
\end{quote}


% ================================================================
\section{Ideja: stroj koji čita riječ}
% ================================================================

Zamislite jednostavan sustav za provjeru PIN-a bankomata. On prima
niz znamenki i na kraju kaže: točno ili pogrešno. U svakom trenutku
pamti samo ``koliko smo znakova dosad prihvatili'' --- nema nikakve
druge memorije. Takav sustav je, u biti, konačni automat.

Formalno, konačni automat je idealizirani matematički stroj koji:
\begin{itemize}
  \item ima konačno mnogo \textbf{stanja} (kao moguće ``faze'' rada),
  \item kreće iz fiksnog \textbf{početnog stanja},
  \item za svaki pročitani znak prelazi u novo stanje prema
        \textbf{funkciji prijelaza},
    \item na kraju prihvaća riječ ako je završio u nekom od
      \textbf{završnih stanja}.
\end{itemize}

Kljucno ogranicenje: nema stoga, nema trake, nema ekstre memorije.
Jedino sto stroj ``zna'' jest u kojem je stanju --- i to je sve.

% ================================================================
\section{Deterministički konačni automat (KA)}
% ================================================================

\begin{definicija}[Konačni automat]
\textbf{Konačni automat} (KA) nad abecedom $\Sig$ je uređena petorka
\[
  M = (Q,\, \Sig,\, \delta,\, q_0,\, F),
\]
gdje je:
\begin{itemize}
  \item $Q$ --- konačni skup \textbf{stanja},
  \item $\Sig$ --- abeceda (konačni skup ulaznih znakova),
  \item $\delta : Q \times \Sig \to Q$ --- \textbf{funkcija prijelaza},
  \item $q_0 \in Q$ --- \textbf{početno stanje},
  \item $F \subseteq Q$ --- skup \textbf{završnih stanja}.
\end{itemize}
\end{definicija}

Funkcija prijelaza $\delta$ je srce automata: za svako stanje i svaki
ulazni znak precizno govori u koje stanje automat prelazi. Nema
nedefiniranih prijelaza, nema izbora --- za svaki par $(q, \alpha)$
postoji točno jedan nasljednik $\delta(q, \alpha)$.

\begin{definicija}[Prihvaćanje riječi]
KA $M = (Q, \Sig, \delta, q_0, F)$ \textbf{prihvaća} riječ
$w = \alpha_1 \alpha_2 \cdots \alpha_n \in \Sig^*$ ako postoji niz stanja
$r_0, r_1, \ldots, r_n \in Q$ takav da:
\begin{itemize}
  \item $r_0 = q_0$ (kreće iz početnog stanja),
  \item $\delta(r_{i-1}, \alpha_i) = r_i$ za sve $1 \leq i \leq n$ (prati prijelaze),
  \item $r_n \in F$ (završava u završnom stanju).
\end{itemize}
Jezik koji $M$ prepoznaje je
\[
  L(M) := \{w \in \Sig^* : M \text{ prihvaća } w\}.
\]
\end{definicija}

\begin{napomena}
Za deterministički KA, niz stanja $r_0, r_1, \ldots, r_n$ je
\emph{jedinstven} za zadane $M$ i $w$ --- nema izbora ni slučajnosti.
Automat uvijek "zna" što radi.
\end{napomena}

\begin{primjer}[KA koji prihvaća riječi s parnim brojem $a$-ova]
Nad $\Sig = \{a, b\}$, konstruirajmo KA za jezik svih riječi s
parnim brojem pojavljivanja slova $a$ (uključujući $0$).

Trebamo pratiti samo parnost broja $a$-ova --- dvije mogucnosti:
``parno'' i ``neparno''. To su nase dvije stanice:
\[
  Q = \{q_{\text{par}},\, q_{\text{nepar}}\},
  \quad q_0 = q_{\text{par}},
  \quad F = \{q_{\text{par}}\}.
\]
Funkcija prijelaza:
\begin{center}
\renewcommand{\arraystretch}{1.4}
\begin{tabular}{c|cc}
  & $a$ & $b$ \\ \hline
  $q_{\text{par}}$   & $q_{\text{nepar}}$ & $q_{\text{par}}$ \\
  $q_{\text{nepar}}$ & $q_{\text{par}}$   & $q_{\text{nepar}}$ \\
\end{tabular}
\end{center}
Riječ $aabba$: $q_{\text{par}} \xrightarrow{a} q_{\text{nepar}}
\xrightarrow{a} q_{\text{par}} \xrightarrow{b} q_{\text{par}}
\xrightarrow{b} q_{\text{par}} \xrightarrow{a} q_{\text{nepar}}$.
Završava u $q_{\text{nepar}} \notin F$ --- riječ nije prihvaćena (ima $3$ slova $a$). \checkmark
\end{primjer}

\begin{primjer}[KA koji prihvaća riječi koje završavaju na $ab$]
Nad $\Sig = \{a, b\}$, želimo prepoznati sve riječi koje završavaju
nizom $ab$. Automat mora pamtiti ``suffix koji smo dosad vidjeli'',
ali samo onoliko koliko je potrebno:
\[
  Q = \{q_0,\, q_a,\, q_{ab}\}, \quad F = \{q_{ab}\}.
\]
\begin{center}
\renewcommand{\arraystretch}{1.4}
\begin{tabular}{c|cc}
  & $a$ & $b$ \\ \hline
  $q_0$  & $q_a$  & $q_0$ \\
  $q_a$  & $q_a$  & $q_{ab}$ \\
  $q_{ab}$ & $q_a$ & $q_0$ \\
\end{tabular}
\end{center}
Stanje $q_a$ znači: ``zadnji pročitani znak je bio $a$''. Stanje $q_{ab}$:
``zadnja dva znaka su bila $ab$''. Iz $q_{ab}$, ako pročitamo $a$, idemo
u $q_a$ (novi potencijalni početak); ako pročitamo $b$, gubimo trag
i vraćamo se u $q_0$.
\end{primjer}

% ================================================================
\section{Nedeterministički konačni automat (NKA)}
% ================================================================

Deterministički KA je konceptualno jasan, ali često nespretan za
dizajnirati. Prirodniji pristup je često \emph{nedeterminističan}:
u svakom trenutku automat smije ``pogoditi'' koji će prijelaz napraviti,
a prihvaća riječ ako \emph{postoji barem jedna grana} koja vodi u
završno stanje.

\begin{definicija}[Nedeterministički KA]
\textbf{Nedeterministički konačni automat} (NKA) nad $\Sig$ je
uređena petorka $(Q, \Sig, \Delta, q_0, F)$, gdje je sve kao kod KA
osim sto je funkcija prijelaza zamijenjena \textbf{relacijom prijelaza}:
\[
  \Delta \subseteq Q \times \Sig_\eps \times Q,
\]
gdje je $\Sig_\eps := \Sig \cup \{\eps\}$. Prijelaz $(q, \alpha, p) \in \Delta$
znači: ``iz stanja $q$, čitajući znak $\alpha$ (ili, ako je $\alpha = \eps$,
bez čitanja), može se preći u stanje $p$''.

NKA prihvaća riječ $w$ ako postoji \emph{barem jedan} način čitanja
$w$ koji vodi iz $q_0$ u neko završno stanje.
\end{definicija}

\begin{kljucno}
Nedeterminizam je matematički trik: NKA ``pogađa'' točnu granu.
U stvarnosti takav stroj ne postoji --- ali matematički je koristan
jer je često mnogo lakši za dizajnirati od ekvivalentnog KA.
Ključni rezultat: svaki NKA ima ekvivalentni KA!
\end{kljucno}

\subsection{Epsilon-prijelazi}

Posebno su korisni \textbf{$\eps$-prijelazi}: automat prelazi iz jednog
stanja u drugo \emph{bez čitanja ijednog znaka}. Ovo se čini kao
varanje, ali zapravo samo kaže: ``automat smije besplatno proći kroz
neka stanja''.

\begin{definicija}[$\eps$-ljuska]
Za skup stanja $S \subseteq Q$, \textbf{$\eps$-ljuska} od $S$ je skup
svih stanja dostupnih iz $S$ iskljucivo $\eps$-prijelazima:
\[
  [S]_\eps := \{q' \in Q : \text{postoji staza } \eps\text{-prijelaza od nekog } q \in S \text{ do } q'\}.
\]
\end{definicija}

\begin{primjer}[$\eps$-prijelazi u praksi]
Kad želimo spojiti dva NKA $N_1$ i $N_2$ tako da prepoznaju
konkatenaciju $L(N_1) L(N_2)$, to radimo tako da iz svakog završnog
stanja od $N_1$ dodamo $\eps$-prijelaz prema početnom stanju od $N_2$.
Niti jedan znak se ne ``gubi'' --- samo automati postaju lijepo spojeni.
\end{primjer}

% ================================================================
\section{Ekvivalentnost KA i NKA}
% ================================================================

\begin{propozicija}[Sipser, teorem 1.39]
Jezik je KA-regularan ako i samo ako je NKA-regularan.
\end{propozicija}

Smjer $(\Rightarrow)$ je trivijalan: svaki KA je specijalan slucaj NKA
(bez nedeterminizma i $\eps$-prijelaza). Smjer $(\Leftarrow)$ je
zanimljiviji i koristi \textbf{partitivnu konstrukciju}.

\subsection{Partitivna konstrukcija}

Ideja: ako NKA u svakom trenutku može biti u više stanja odjednom,
zasto ne bismo cijeli skup mogućih stanja proglasili jednim stanjem
novog, determinističkog automata?

\begin{definicija}[Partitivna konstrukcija]
Neka je $N = (Q, \Sig, \Delta, q_0, F)$ NKA. Definiramo KA
\[
  P(N) = \bigl(\mathcal{P}(Q),\; \Sig,\; \delta,\; [\{q_0\}]_\eps,\;
  \{T \subseteq Q : T \cap F \neq \emptyset\}\bigr),
\]
gdje je
\[
  \delta(T, \alpha) := \bigl[\{q \in Q : (\exists p \in T)\; (p, \alpha, q) \in \Delta\}\bigr]_\eps.
\]
\end{definicija}

Intuicija: stanje $P(N)$-a je skup stanja u kojima $N$ \emph{može}
biti. Kad pročitamo znak $\alpha$, uzimamo sve $\Delta$-nasljednike
svakog stanja iz skupa, a zatim zatvorimo na $\eps$-prijelaze.

\begin{napomena}
Ako $N$ ima $n$ stanja, $P(N)$ ima $2^n$ stanja --- eksponencijalan
rast! U praksi se većinom koriste samo dostupna stanja (ona do kojih
se uopće može doći iz početnog), kojih je obično znatno manje.
\end{napomena}

\begin{primjer}[Partitivna konstrukcija]
Neka je $N$ NKA s $Q = \{1, 2, 3\}$, $q_0 = 1$, $F = \{3\}$ i
prijelazima:
\[
  \Delta = \{(1,a,1),\; (1,a,2),\; (2,b,3),\; (3,b,3)\}.
\]
Ovaj NKA prepoznaje rijeci oblika $a^n b^k$ za $n \geq 1$, $k \geq 1$
(nejasno, ali pogledajmo partitivnu konstrukciju):

Početno stanje $P(N)$-a je $\{1\}$. Pratimo:
\begin{itemize}
  \item $\delta(\{1\}, a) = \{1, 2\}$ (jer $(1,a,1)$ i $(1,a,2)$)
  \item $\delta(\{1,2\}, a) = \{1, 2\}$ (iz 1 kao gore, iz 2 nema)
  \item $\delta(\{1,2\}, b) = \{3\}$ (samo $(2,b,3)$)
  \item $\delta(\{3\}, b) = \{3\}$ i $\delta(\{3\}, a) = \emptyset$
\end{itemize}
Dostupna stanja: $\{1\}$, $\{1,2\}$, $\{3\}$, $\emptyset$ ---
samo $4$ stanja umjesto $2^3 = 8$.
\end{primjer}

% ================================================================
\section{Zatvorenost na skupovne i jezicne operacije}
% ================================================================

Jedna od ključnih osobina regularnih jezika jest da su zatvoreni na
mnoge prirodne operacije.

\subsection{Skupovne operacije}

\begin{propozicija}
Klasa KA-regularnih jezika zatvorena je na \textbf{uniju}, \textbf{presjek}
i \textbf{komplement}.
\end{propozicija}

Dokaz za uniju i presjek koristi \textbf{Kartezijevu konstrukciju}:
ako $M_1$ prepoznaje $L_1$ i $M_2$ prepoznaje $L_2$, konstruiramo
KA koji ``istovremeno'' simulira oba:
\[
  M_1 \times M_2 := (Q_1 \times Q_2,\; \Sig,\; \delta,\; (q_1, q_2),\; F),
\]
gdje $\delta((q,q'), \alpha) := (\delta_1(q,\alpha), \delta_2(q',\alpha))$.

Za \textbf{uniju} stavimo $F = Q_1 \times F_2 \cup F_1 \times Q_2$.
Za \textbf{presjek} stavimo $F = F_1 \times F_2$.
Za \textbf{komplement} samo zamijenimo završna s nezavršnima: $F' = Q \setminus F$.

\subsection{Jezicne operacije}

\begin{propozicija}
Klasa NKA-regularnih (= regularnih) jezika zatvorena je na
\textbf{konkatenaciju}, \textbf{Kleenevu zvijezdu} i \textbf{Kleenejev plus}.
\end{propozicija}

Ove konstrukcije su elegantnije s NKA:
\begin{itemize}
    \item \textbf{Unija $L \cup L'$}: dodamo novo početno stanje s
      $\eps$-prijelazima prema početnim stanjima oba automata.
    \item \textbf{Konkatenacija $LL'$}: iz svakog završnog stanja
      od $N$ dodamo $\eps$-prijelaz prema početnom stanju od $N'$.
      Stara završna stanja od $N$ više nisu završna.
    \item \textbf{Kleeneva zvijezda $L^*$}: dodamo novo početno završno
      stanje s $\eps$-prijelazima prema starim početnim stanjima,
      te iz svakog završnog stanja $\eps$-prijelaz natrag prema
      početnom stanju.
\end{itemize}

% ================================================================
\section{Kako dizajnirati KA: heuristike}
% ================================================================

Dizajn KA je vjestina koja se uci vjezbanjem. Nekoliko heuristika:

\begin{kljucno}
\textbf{Pitajte se: sto KA mora pamtiti?}

Stanja KA su različite ``situacije'' u kojima se automat može naći.
Pitanje je: koje informacije o dosad pročitanom prefiksu su relevantne
za odluku o prihvaćanju?

Primjeri:
\begin{itemize}
  \item ``Zadnji pročitani znak'' --- $|\Sig|$ stanja.
  \item ``Parnost broja $a$-ova'' --- 2 stanja.
  \item ``Ostatak duljine modulo $k$'' --- $k$ stanja.
  \item ``Prefiks duljine do $k$'' --- konačno mnogo stanja.
\end{itemize}
Ako trebate pamtiti nešto \emph{neograničeno} (npr.\ točni broj $a$-ova)
--- KA za to ne postoji.
\end{kljucno}

\begin{zadatak}
\begin{enumerate}[label=(\alph*), itemsep=4pt]
  \item Konstruirajte KA nad $\{a,b\}$ koji prepoznaje jezik svih rijeci
        cija duljina je djeljiva s $3$.
  \item Konstruirajte NKA koji prepoznaje jezik svih rijeci nad $\{0,1\}$
        koje završavaju s $010$.
  \item Primijenite partitivnu konstrukciju na NKA iz (b) da dobijete
        ekvivalentni KA.
  \item Konstruirajte KA koji prepoznaje komplement jezika iz (a).
\end{enumerate}
\end{zadatak}

% ================================================================
\section*{Sazetak}
% ================================================================

Konačni automat čita riječ znak po znak, ne može se vraćati unatrag,
i pamti samo stanje u kojemu je. Usprkos ovim jakim ograničenjima,
konačni automati prepoznaju točno regularnu klasu jezika.

\medskip
\noindent
\begin{tabular}{ll}
  $M = (Q, \Sig, \delta, q_0, F)$ & deterministički KA \\
  $N = (Q, \Sig, \Delta, q_0, F)$ & nedeterministički NKA \\
  $[S]_\eps$ & $\eps$-ljuska skupa stanja \\
  $P(N)$ & partitivna konstrukcija (NKA $\to$ KA) \\
  $M_1 \times M_2$ & Kartezijeva konstrukcija (unija/presjek) \\
\end{tabular}

\medskip
\noindent U sljedećem poglavlju zatvaramo krug: pokazujemo da su
KA, NKA, regularni izrazi i linearne gramatike sve \emph{isti} formalizam,
samo u različitom ruhu.

